%---------------------------------------------------------------
\chapter{Implementace}
\label{ch:implementace}
%---------------------------------------------------------------

Tato kapitola popisuje implementaci jednotlivých komponent systému SwimAth podle návrhu z~kapitoly~\ref{ch:navrh}.
Kapitola je strukturována od~předzpracování videa přes ML pipeline a~webovou aplikaci až~po~testovací strategii a~nasazení.

\section{Předzpracování videa}
\label{sec:predzpracovani}

% TODO: OpenCV pipeline
% - VideoCapture wrapper: otevření videa, kontrola codec/fps/resolution
% - Validace formátu: MP4/MOV, H.264/H.265 codec
% - Normalizace: resize na max 1080p (zachování aspect ratio), normalizace FPS na 30
% - FFmpeg integrace: konverze nestandardních formátů
% - Frame extraction: sekvenční čtení framů do numpy array
% - Buffering strategie: batch processing po N framech (memory efficient)

\section{ML pipeline}
\label{sec:ml_pipeline}

% TODO: Implementace ML pipeline jako SW komponenty

\subsection{ONNX Runtime inference}

% TODO: Popis inference pipeline
% - Načtení modelu: ONNX Runtime session s GPU ExecutionProvider
% - Batched inference: zpracování N framů najednou (optimalizace GPU utilization)
% - Pre-processing: resize, normalizace, top-down crop podle bounding boxu
% - Post-processing: heatmap → keypoints (argmax + subpixel refinement)
% - Fallback: pokud GPU není dostupný, CPU inference (pomalejší, ale funkční)

\subsection{Strategy pattern pro PE modely}

% TODO: Popis implementace Strategy pattern
% - PoseEstimatorBase (ABC): predict(frames) → keypoints
% - MediaPipePoseEstimator: wrapper pro MediaPipe (baseline)
% - RTMPosePoseEstimator: ONNX inference RTMPose-m
% - ViTPosePoseEstimator: ONNX inference ViTPose-B
% - Výběr strategie na základě konfigurace (env variable / admin setting)
% - Konfigurační rozhraní: model_path, input_size, keypoint_count, confidence_threshold

\subsection{Pipeline orchestrace}

% TODO: Popis pipeline flow
% - Pipeline třída: orchestrace kroků (PE → ragdoll → metrics → DTW → feedback)
% - Každý krok jako samostatná třída s rozhraním process(input) → output
% - Error handling: selhání na jednom framu nezastaví celou analýzu
% - Progress callback: po každém kroku update progressu přes Redis

\section{Backend}
\label{sec:backend}

% TODO: Implementace FastAPI backendu

\subsection{Struktura projektu}

% TODO: Diagram adresářové struktury (listing)
% swimath-api/
% ├── app/
% │   ├── main.py              # FastAPI app, CORS, lifecycle
% │   ├── config.py            # Pydantic BaseSettings
% │   ├── api/                 # Endpointy (upload, analysis, history, websocket)
% │   ├── ml/                  # ML pipeline (PE, biomechanics, comparator, feedback)
% │   ├── processing/          # Video zpracování (handler, extractor, pipeline)
% │   ├── references/          # Referenční databáze (manager, templates)
% │   ├── db/                  # SQLAlchemy modely, Alembic migrace
% │   └── utils/               # Pomocné funkce

\subsection{Routing a dependency injection}

% TODO: FastAPI specifika
% - APIRouter pro modularitu endpointů
% - Dependency injection: DB session, Celery app, config
% - Pydantic modely pro validaci request/response
% - Middleware: CORS, request logging, error handling

\subsection{Pydantic modely}

% TODO: Ukázkové modely (listing)
% - UploadParams: skill_level (Enum), camera_view (Enum), swim_style (Enum)
% - AnalysisResponse: id, status, metrics, deviations, feedback, overall_score
% - ProgressMessage: progress (int), stage (str), message (str)

\section{Frontend}
\label{sec:frontend}

% TODO: Implementace Vue.js 3 frontendu

\subsection{Vue.js 3 Composition API}

% TODO: Architektura frontendu
% - Composition API s <script setup> syntaxí
% - Pinia store pro globální stav (aktuální analýza, historie)
% - Vue Router pro navigaci (upload → progress → results → history)
% - Axios instance s interceptory pro error handling

\subsection{Upload komponenta}

% TODO: Popis upload flow
% - Drag & drop + file input
% - Klientská validace (formát, velikost)
% - Formulář pro výběr parametrů (skill_level, camera_view, swim_style)
% - Axios POST s FormData
% - Redirect na progress view

\subsection{WebSocket progress}

% TODO: Popis real-time progress
% - WebSocket spojení na /ws/progress/{task_id}
% - Reactive progress bar (Vue reactive refs)
% - Stage indikátor (jaký krok pipeline právě běží)
% - Auto-redirect na výsledky po dokončení

\subsection{Video overlay --- Canvas API}

% TODO: Popis vizualizace pózy
% - HTML5 Canvas overlay přes <video> element
% - Vykreslení skeleton (keypoints + connections)
% - Synchronizace canvas s video playback (requestAnimationFrame)
% - Barevné kódování: zelená (OK), žlutá (minor), červená (severe) pro problematické klouby
% - Timeline vizualizace záběrových cyklů

\section{Asynchronní zpracování}
\label{sec:async_impl}

% TODO: Implementace Celery pipeline

\subsection{Celery workers}

% TODO: Konfigurace a implementace
% - Celery app s Redis brokerem
% - Task: analyze_video(analysis_id) — hlavní processing task
% - Konfigurace: concurrency, prefetch multiplier, task time limit
% - GPU worker: dedikovaný worker s GPU přístupem pro ML inference
% - CPU worker: pro lehké úlohy (validace, export)

\subsection{WebSocket progress bridge}

% TODO: Redis pub/sub → WebSocket
% - Worker publikuje progress do Redis kanálu
% - FastAPI WebSocket endpoint odebírá z Redis kanálu
% - Bridge zajišťuje real-time doručení progress updates klientovi

\subsection{Task state management}

% TODO: Stavový diagram tasku
% - Stavy: PENDING, STARTED, PROGRESS, SUCCESS, FAILURE, REVOKED
% - Retry logika: max 3 retries, exponential backoff
% - Timeout: task time limit pro prevenci zombie tasků
% - Cleanup: automatické smazání starých tasků

\section{Databáze}
\label{sec:databaze}

% TODO: Implementace persistence vrstvy

\subsection{SQLAlchemy ORM}

% TODO: Popis ORM modelů
% - Deklarativní mapování (DeclarativeBase)
% - Modely: Analysis, Metric, Deviation, Feedback, ReferenceTemplate
% - Vztahy: relationship() s lazy loading
% - JSON column pro flexibilní metadata

\subsection{Alembic migrace}

% TODO: Migrace workflow
% - alembic revision --autogenerate pro generování migrací
% - alembic upgrade head při startu aplikace
% - Verzování schématu v git

\subsection{Connection pooling}

% TODO: SQLAlchemy connection pool
% - Pool size, max overflow, pool recycle
% - Async session (pokud FastAPI async endpoints)

\section{Referenční databáze}
\label{sec:referencni_db}

% TODO: Implementace referenčního systému

% - Struktura: skill_level / swim_style / camera_view / template.json
% - JSON schéma: temporální sekvence keypointů (DBA výstup) + variance envelope
% - Generování: DBA iterativní algoritmus na trénovacích záběrech ze SwimXYZ
% - Verzování šablon: verze v JSON metadata, migrace při změně formátu
% - Lazy loading: šablony se načítají při prvním požadavku a cacheují v paměti

\section{Testovací strategie}
\label{sec:testovani}

% TODO: Detailní popis testovacího přístupu

\subsection{Testovací pyramida}

% TODO: Diagram testovací pyramidy (figure)
% - Unit testy (základ): jednotlivé funkce, třídy, moduly
% - Integrační testy (střed): API endpointy, DB operace, Celery tasky
% - E2E testy (špička): celý upload → výsledky flow

\subsection{Unit testy}

% TODO: pytest specifika
% - Fixtures pro DB session, test client, mock ML models
% - Parametrizované testy pro různé kombinace (styl × pohled × úroveň)
% - Mock PE modelu pro rychlé testy (předpočítané keypoints)
% - Coverage target: >80%

\subsection{Integrační testy}

% TODO: FastAPI TestClient
% - Test upload endpointu (validace, vytvoření záznamu)
% - Test analysis endpointu (vrácení výsledků)
% - Test WebSocket (progress updates)
% - Test Celery tasku s mock ML modelem

\subsection{E2E testy}

% TODO: Playwright testy
% - Upload flow: nahrání videa → progress → výsledky
% - Video overlay: kontrola renderování
% - Historie: zobrazení předchozích analýz

\section{CI/CD a deployment}
\label{sec:cicd}

% TODO: Popis nasazení a automatizace

\subsection{Docker}

% TODO: Multi-stage Dockerfile
% - Build stage: instalace závislostí, kompilace
% - Runtime stage: minimální image, non-root user
% - docker-compose: frontend, api, worker, db, redis
% - GPU podpora: nvidia-docker pro ML worker

\subsection{CI pipeline}

% TODO: GitHub Actions workflow
% - Lint: ruff, mypy
% - Test: pytest s coverage reportem
% - Build: Docker image
% - (Budoucí) Deploy: push to registry, deploy to server
